\subsectionaddtolist{استخراج فرمول‌های بازیابی}

{آ) فرمول صریح برای بازیابی $k$:}
با توجه به اینکه دو امضا با نانس یکسان $k$ و مقدار یکسان $r$ تولید شده‌اند، معادله‌های امضا به صورت زیر هستند:
\[ s_1 \equiv k^{-1}(H(M_1) + xr) \pmod q \]
\[ s_2 \equiv k^{-1}(H(M_2) + xr) \pmod q \]

با تفاضل دو معادله خواهیم داشت:
\[ s_1 - s_2 \equiv k^{-1}(H(M_1) - H(M_2)) \pmod q \]

در نتیجه، فرمول صریح برای بازیابی $k$ برابر است با:
\[ k \equiv (s_1 - s_2)^{-1}(H(M_1) - H(M_2)) \pmod q \]

{ب) فرمول صریح برای بازیابی کلید خصوصی $x$:}
با استفاده از معادله امضای اول:
\[ s_1 \equiv k^{-1}(H(M_1) + xr) \pmod q \implies k \cdot s_1 \equiv H(M_1) + xr \pmod q \]
\[ xr \equiv k \cdot s_1 - H(M_1) \pmod q \]

با ضرب دو طرف در معکوس ضربی $r$ مرتبه $q$ داریم:
\[ x \equiv r^{-1}(k \cdot s_1 - H(M_1)) \pmod q \]

\subsectionaddtolist{محاسبه کلید خصوصی}

پارامترهای داده شده:
\[ p = 23, \quad q = 11, \quad g = 2, \quad y = 18 \]
\[ H(M_1) = 4, \quad (r, s_1) = (8, 10) \]
\[ H(M_2) = 5, \quad (r, s_2) = (8, 3) \]

{آ) محاسبه مقدار $k$:}
ابتدا تفاضل‌ها را در پیمانه $q = 11$ محاسبه می‌کنیم:
\[ s_1 - s_2 = 10 - 3 = 7 \pmod{11} \]
\[ H(M_1) - H(M_2) = 4 - 5 = -1 \equiv 10 \pmod{11} \]

معکوس ضربی $7$ در پیمانه $11$ برابر است با $8$ (زیرا $7 \times 8 = 56 \equiv 1 \pmod{11}$).\\
بنابراین:
\[ k \equiv 7^{-1} \times 10 \equiv 8 \times 10 = 80 \equiv 3 \pmod{11} \]
پس مقدار نانس برابر است با: {$k = 3$}

{ب) محاسبه کلید خصوصی آلیس $x$:}
ابتدا معکوس ضربی $r = 8$ را در پیمانه $11$ پیدا می‌کنیم که برابر $7$ است (زیرا $8 \times 7 = 56 \equiv 1 \pmod{11}$).\\
حال مقدار داخل پرانتز را محاسبه می‌کنیم:
\[ k \cdot s_1 - H(M_1) = 3 \times 10 - 4 = 26 \equiv 4 \pmod{11} \]

بنابراین کلید خصوصی $x$ برابر است با:
\[ x \equiv 7 \times 4 = 28 \equiv 6 \pmod{11} \]
پس کلید خصوصی برابر است با: \textbf{$x = 6$}

{ج) تایید صحت محاسبات با تولید کلید عمومی:}
رابطه کلید عمومی در \lr{DSS} به صورت $y = g^x \pmod p$ است:
\[ y = 2^6 \pmod{23} = 64 \pmod{23} \]
از آنجا که $64 = 2 \times 23 + 18$ است، داریم:
\[ y = 18 \]
این مقدار دقیقاً با کلید عمومی داده شده در صورت سوال ($y = 18$) مطابقت دارد، پس محاسبات درست است.

\subsectionaddtolist{مکانیزم استاندارد جلوگیری از استفاده مجدد \lr{Nonce}}

{آ) شناسایی مکانیزم استاندارد:}
مکانیزم استاندارد صنعتی برای تولید دترمینستیک (یکتا و معین) نانس در امضاهای \lr{DSA} و \lr{ECDSA} در سند \textbf{\lr{RFC 6979}} تعریف شده است.

{ب) نحوه تولید نانس در این مکانیزم:}
در این روش، مقدار $k$ به صورت کاملاً دترمینستیک و بر اساس ترکیب کلید خصوصی ($x$) و چکیده پیام ($H(M)$) با استفاده از توابع هش رمزنگاری مانند \lr{HMAC-SHA256} تولید می‌شود. به این ترتیب، برای یک پیام و کلید خصوصی مشخص، همواره یک $k$ یکتا و کاملاً تصادفی‌نما تولید خواهد شد.

{ج) دلیل از بین رفتن خطرات مرتبط با تولیدکننده‌های اعداد تصادفی ضعیف:}
از آنجا که فرآیند تولید $k$ دیگر به هیچ منبع سخت‌افزاری یا نرم‌افزاری تولید اعداد تصادفی سیستم در لحظه امضا متکی نیست، خطرات ناشی از آنتروپی ضعیف، نقص‌های فنی یا حملات بهره‌برداری از تولیدکننده تصادفی (\lr{RNG}) به طور کامل حذف می‌شود. اگر یک پیام دو بار امضا شود، نانس یکسانی تولید شده و امضای کاملاً مشابهی ایجاد می‌کند که هیچ اطلاعاتی در مورد کلید خصوصی فاش نخواهد کرد.